\begin{problem}[2020第37届全国中学生物理竞赛复赛][星际飞行器相对论通信与货物传输]{121}
    星际飞行器甲、乙、丙在惯性系 $\Sigma$（坐标系 $O-xy$）中观测到飞行器甲和乙在同一直线上朝 $y$ 轴负方向匀速飞行，飞行器丙在另一直线上朝 $y$ 轴正方向匀速飞行；两条飞行直线相互平行，相距 $d$。三艘飞行器的速度在惯性系 $\Sigma$ 中分别为 $\mathbf{v}_{\mathrm{甲}\Sigma} = (-2c/3) \hat{\mathbf{y}}$，$\mathbf{v}_{\mathrm{乙}\Sigma} = (-c/3) \hat{\mathbf{y}}$，$\mathbf{v}_{\mathrm{丙}\Sigma} = (2c/3) \hat{\mathbf{y}}$（$\hat{\mathbf{y}}$ 是沿 $y$ 轴正方向的单位矢量），图中飞行器旁的箭头表示其飞行速度的方向。假设飞行器甲、乙、丙的尺寸远小于 $d$。
    \begin{subproblem}
        某时刻，飞行器甲向乙发射一个光信号，乙收到该信号后立即将其反射回甲。据甲上的原子钟读数，该光信号从发出到返回共经历了时间 $(\Delta t_{\text{光信号}})_{\mathrm{甲}} = T$。试求分别从惯性系 $\Sigma$ 和乙上的观测者来看，光信号从甲发出到返回甲，所经历的时间各是多少？从甲上的观测者来看，从甲收到返回的光信号到甲追上乙需要多少时间？
    \end{subproblem}
    \begin{subproblem}
        当飞行器甲和丙在惯性系 $\Sigma$ 中相距最近时，从甲发射一个小货物（质量远小于飞行器的质量，尺寸可忽略），小货物在惯性系 $\Sigma$ 中的速度大小为 $v_{\mathrm{货}\Sigma} = 3c/4$。为了让丙能接收到该货物，从甲上的观测者来看，该货物发射速度的大小和方向是多少？从乙上的观测者来看，货物从甲发射到被丙接收所需时间 $(\Delta t_{\text{货}})_{\mathrm{乙}}$ 是多少？
    \end{subproblem}
\end{problem}